\scanpage{48}
l'ouvert cherché [N.B. la locale liberté de $M$ en
$x$ est en effet une condition \uline{ponctuelle}
et non seulement locale, \deleted{\unsure{c'est|exacte}} compatible avec
l'extension des corps de base $k(x)$ \ldots]

\noindent\rule{3cm}{0.4pt}\\[-0.8ex]
\rule{2.8cm}{0.4pt}

Sous les réserves habituelles, \uline{si} $g:X_{/T}\unsure{\mathbin{-}|\to}Y_{/T}$ \uline{est de
type fini}, \uline{alors} $f_\ast(g)$ \uline{est de type fini}.
[ $S$ \deleted{\unsure{est}} \uline{loc. noeth.} )

On est ramené au cas où $Y=T$. \deleted{$X$ affine
sur $T$. On peut \unsure{se supposer}} \inserted{Par}
hypothèse, \deleted{\unsure{toute|$X$}} multisection d'un
$X_s$ sur un $T_s$ est contenue dans un ouvert \inserted{affine $U$}\deleted{\unsure{aff.|$X$}}.
\unsure{Prouvons|Prenons}, qu'\unsure{une st|une qté|un nb} finie de tels ouverts
affines \deleted{\unsure{$X$}} suffisent. \deleted{\unsure{On voit que ça|Car}
suffira, dans le \unsure{cas}} \inserted{Soit en effet}
% slanted
\margin{\uline{dégager un lemme}}
\[
\deleted{\unsure{X}}\unsure{X^{\nu}|X^{v}}=
\underbrace{X\times_S\mathbin{-}\times_S X}_{\nu}\,,
\]
pour tout ouvert affine
$U$ de $X$, soit $U^\nu=\underbrace{U\times_S\mathbin{-}\times_S U}$, $U^\nu$ est un
ouvert de $X^\nu$. \deleted{\unsure{On verra que|Par}} \inserted{Comme $X^\nu$ est}
noethérien, il existe une \uline{st finie} d'ouverts
affines $U_i$ tels que $\bigcup U_i^\nu=\bigcup U^\nu$.
Soient $(x_1,\mathbin{-},x_\nu)$\deleted{$\in\unsure{\lvert X\rvert^\nu}$ \unsure{que} \unsure{les éléments d'un même ouvert affine}
$U$. Alors} soit $x\in X^{\deleted{\unsure{n}}\nu}$ de projections les $x_i$,
Alors les $x_i$ sont dans un ouvert affine
% archive sheet number: 47
